\documentclass[11pt]{article}

\usepackage[utf8]{inputenc}
\usepackage[T1]{fontenc}
\usepackage{amsmath,amssymb,amsfonts}
\usepackage{bm}

\title{MacroTaylor: Rank-2 Laplace Approximation for the Posterior over $p$ \\
       \large (Moffatt 2007 scaling convention)}
\author{}
\date{}

\begin{document}

\maketitle

This note is a rewrite of \texttt{MacroTaylor\_2\_chatgpt\_reseach.tex} in the
\emph{Moffatt (2007) scaling convention} for the prior and posterior
covariances.  The mathematics is identical; the only change is the
choice of which object carries the explicit $1/N_{\text{ch}}$ factor.

\paragraph{Scaling convention.}
Following Moffatt (2007), we represent the prior and posterior
covariances over $p$ by their \emph{per-channel} (intensive) scaled
counterparts $\Sigma$ and $\Sigma_{\text{post}}$, which are independent
of the channel count $N_{\text{ch}}$.  The \emph{actual} prior and
posterior covariances carry an explicit $1/N_{\text{ch}}$:
\[
  p \sim \mathcal{N}\!\bigl(\mu,\ \tfrac{1}{N_{\text{ch}}}\Sigma\bigr),
  \qquad
  p\mid y
  \approx
  \mathcal{N}\!\bigl(p^\star,\ \tfrac{1}{N_{\text{ch}}}\Sigma_{\text{post}}\bigr).
\]
Adding channels makes the actual covariance shrink as $1/N_{\text{ch}}$,
while $\Sigma$ stays fixed.  This is the natural convention for an
ensemble of $N_{\text{ch}}$ i.i.d.\ channels and reduces cleanly to the
single-channel formulas of Moffatt (2007) when $\sigma^2 = 0$.

The \textbf{main scheme} is the \emph{rank-2} Laplace approximation
based on the \emph{exact} Gaussian likelihood (including the
$\tfrac12\log V$ term).  As a \emph{pedagogical limit} and for
connection to EKF-style approximations, we also present a
\emph{rank-1} ``quasi-likelihood'' scheme in which the $\log V$ term is
dropped.

The goals are:
\begin{itemize}
  \item to define the exact Gaussian likelihood and Moffatt-scaled
        prior in the macroscopic coordinates $p$;
  \item to derive the exact gradient and Hessian of the negative
        log-posterior, and show that the Hessian is a rank-2 update
        of the (per-channel) prior precision $\Sigma^{-1}$;
  \item to obtain a closed-form expression for the \emph{scaled
        Laplace covariance} $\Sigma_{\text{post}}$ via a rank-2
        Woodbury inversion;
  \item to derive a mean-update (MAP) formula in which the
        $1/N_{\text{ch}}$ factors cancel and no explicit $\Sigma^{-1}$
        survives;
  \item to show how the rank-1 (no $\log V$) scheme appears as a
        limiting case and how it connects to quasi-likelihood/EKF
        practice.
\end{itemize}

\paragraph{Coordinates, simplex and singular covariance.}
The macroscopic occupancy vector $p$ lives on the simplex
\[
  \mathcal{S} = \Bigl\{p \in \mathbb{R}^{1\times d} :
    p_i \ge 0,\ \sum_i p_i = 1\Bigr\},
\]
a $(d-1)$-dimensional manifold embedded in $\mathbb{R}^{1\times d}$.
The scaled prior covariance $\Sigma$ is therefore singular as a
$d\times d$ matrix in ambient coordinates.  Formally, one may view
$\Sigma^{-1}$ as the inverse on the intrinsic $(d-1)$-dimensional
subspace, or as a pseudo-inverse.  In what follows:
\begin{itemize}
  \item we freely use $\Sigma^{-1}$ \emph{algebraically}
        when applying Sherman--Morrison--Woodbury;
  \item all \emph{implementable} expressions for the mean and
        covariance are written purely in terms of $\Sigma$ (no
        explicit inverse);
  \item the log-determinant of $\Sigma$ (or its pseudo-determinant)
        is a constant in $p$ and does not affect the MAP.
\end{itemize}

\section*{0. Setup and Notation}

We consider:
\begin{itemize}
  \item Macroscopic occupancy vector
    \[
      p \in \mathbb{R}^{1\times d}
    \]
    (row vector of state probabilities).
  \item Per-state conductance (or current) vector
    \[
      \gamma \in \mathbb{R}^{d\times 1}
    \]
    (column vector).
  \item Per-state intrinsic-variance vector
    \[
      \sigma^2 \in \mathbb{R}^{d\times 1}.
    \]
\end{itemize}

The prior over $p$ (restricted to the simplex) is Gaussian in
Moffatt-scaled form:
\[
  p \sim \mathcal{N}\!\bigl(\mu,\ \tfrac{1}{N_{\text{ch}}}\Sigma\bigr),
\]
where $\mu \in \mathbb{R}^{1\times d}$ is the prior mean and
$\Sigma \in \mathbb{R}^{d\times d}$ is the \emph{scaled} (per-channel)
prior covariance.  The corresponding prior precision is
\[
  \Lambda_p
  = \bigl(\tfrac{1}{N_{\text{ch}}}\Sigma\bigr)^{-1}
  = N_{\text{ch}}\,\Sigma^{-1}.
\]

Over an interval we observe a scalar averaged current $y$ which, for a
given $p$, we model as Gaussian:
\[
  y \mid p
  \sim
  \mathcal{N}\bigl(\mu_y(p), V(p)\bigr),
\]
with
\[
  \mu_y(p) = N_{\text{ch}}(p\cdot\gamma),
  \qquad
  V(p) = \epsilon^2 + N_{\text{ch}}(p\cdot\sigma^2),
\]
where $\epsilon^2$ is the extrinsic (instrument) variance and
$N_{\text{ch}}(p\cdot\sigma^2)$ is the intrinsic interval variance
from channel stochasticity.

Define the residual
\[
  \delta(p) = y - \mu_y(p) = y - N_{\text{ch}}(p\cdot\gamma).
\]

We will also use the \emph{tilted} vector
\[
  v(p)
  =
  \gamma + \frac{\delta(p)}{V(p)}\,\sigma^2
  \in \mathbb{R}^{d\times 1},
\]
which will factor the Hessian.

\section*{1. Exact Gaussian Likelihood and Moffatt-scaled Prior}

We start from the standard Gaussian forms and derive the exact
negative log-posterior.  We then introduce an ``energy'' that is
twice the negative log-posterior, which simplifies later algebra.

\subsection*{1.1 Exact likelihood}

Given $p$, the likelihood is
\[
  p(y\mid p)
  =
  \frac{1}{\sqrt{2\pi V(p)}}
  \exp\!\left(
    -\frac{1}{2}\frac{\delta(p)^2}{V(p)}
  \right).
\]

The log-likelihood is
\[
  \log p(y\mid p)
  =
  -\frac{1}{2}\log(2\pi)
  -\frac{1}{2}\log V(p)
  -\frac{1}{2}\frac{\delta(p)^2}{V(p)}.
\]

Up to an additive constant $-\frac12\log(2\pi)$, the \emph{negative}
log-likelihood is
\[
  F_{\text{like}}(p)
  =
  \frac{1}{2}\log V(p)
  + \frac{1}{2}\frac{\delta(p)^2}{V(p)}
  \quad(+\ \text{const}).
\]

\subsection*{1.2 Moffatt-scaled prior}

The Moffatt-scaled prior density is
\[
  p(p)
  \propto
  \exp\!\left(
    -\frac{1}{2}(p-\mu)\,\bigl(N_{\text{ch}}\Sigma^{-1}\bigr)\,(p-\mu)^\top
  \right),
\]
since the actual prior covariance is $\tfrac{1}{N_{\text{ch}}}\Sigma$
and its precision is $N_{\text{ch}}\Sigma^{-1}$.  The corresponding
negative log-prior (up to a constant) is
\[
  F_{\text{prior}}(p)
  =
  \frac{N_{\text{ch}}}{2}(p-\mu)\,\Sigma^{-1}\,(p-\mu)^\top
  \quad(+\ \text{const}).
\]

The log-determinant term $\frac12\log\det\bigl(\tfrac{1}{N_{\text{ch}}}\Sigma\bigr)$
is a constant in $p$ and therefore does not affect the MAP; we drop
it throughout.  However, we should note that the posterior covariance
does depend on $p$, so this is actually a simplification to make the
calculations manageable.

\subsection*{1.3 Exact negative log-posterior and energy}

The (unnormalised) posterior is
\[
  p(p\mid y) \propto p(y\mid p)\,p(p),
\]
so the exact negative log-posterior (up to a constant) is
\[
  F_{\text{exact}}(p)
  =
  F_{\text{like}}(p) + F_{\text{prior}}(p).
\]

Explicitly:
\[
  F_{\text{exact}}(p)
  =
  \frac{1}{2}\log V(p)
  + \frac{1}{2}\frac{\delta(p)^2}{V(p)}
  + \frac{N_{\text{ch}}}{2}(p-\mu)\,\Sigma^{-1}\,(p-\mu)^\top
  \quad(+\ \text{const}).
\]

We will minimise $F_{\text{exact}}(p)$ to obtain the exact MAP.  For
algebraic convenience, define the associated ``energy''
\[
  E_{\text{exact}}(p)
  := 2F_{\text{exact}}(p).
\]

Up to an additive constant,
\[
  E_{\text{exact}}(p)
  =
  \underbrace{\log V(p) + \frac{\delta(p)^2}{V(p)}}_{L(p) + M(p)}
  +
  \underbrace{N_{\text{ch}}\,(p-\mu)\,\Sigma^{-1}\,(p-\mu)^\top}_{P(p)}.
\]

Thus:
\[
  E_{\text{exact}}(p) = L(p) + M(p) + P(p).
\]

Multiplying by a positive constant does not change the MAP; moreover,
if we write $\mathbf{H}(p) = \nabla^2 E_{\text{exact}}(p)$, then the
\emph{actual} Laplace covariance of the posterior is
\[
  \tfrac{1}{N_{\text{ch}}}\Sigma_{\text{Lap}}
  \approx
  \bigl[\nabla^2 F_{\text{exact}}(p^\star)\bigr]^{-1}
  =
  2\,\mathbf{H}(p^\star)^{-1},
\]
where $p^\star$ is the MAP and $\Sigma_{\text{Lap}}$ denotes the
\emph{scaled} (per-channel) Laplace covariance.

\section*{2. Gradient and Hessian of the Exact Energy}

We now derive the gradient and Hessian of $E_{\text{exact}}(p)$.  We
keep track of three contributions:
\[
  L(p) = \log V(p),
  \quad
  M(p) = \frac{\delta(p)^2}{V(p)},
  \quad
  P(p) = N_{\text{ch}}\,(p-\mu)\,\Sigma^{-1}\,(p-\mu)^\top.
\]

Here $p$ is a row vector, $\gamma$ and $\sigma^2$ are column vectors.

Recall:
\[
  \delta(p) = y - N_{\text{ch}}(p\cdot\gamma),
  \qquad
  V(p) = \epsilon^2 + N_{\text{ch}}(p\cdot\sigma^2).
\]

Then
\[
  \nabla\delta(p) = -N_{\text{ch}}\gamma^\top,
  \qquad
  \nabla V(p) = N_{\text{ch}}(\sigma^2)^\top,
\]
both row vectors, and $\nabla^2\delta = 0$, $\nabla^2 V = 0$.

\subsection*{2.1 Prior gradient and Hessian}

Write
\[
  P(p) = (p-\mu)\,\Lambda_p\,(p-\mu)^\top,
  \qquad
  \Lambda_p = N_{\text{ch}}\Sigma^{-1}.
\]

Since $\Lambda_p$ is symmetric,
\[
  \nabla P(p) = 2(p-\mu)\Lambda_p
              = 2N_{\text{ch}}(p-\mu)\Sigma^{-1},
\]
and
\[
  \nabla^2 P(p) = 2\Lambda_p = 2N_{\text{ch}}\Sigma^{-1}.
\]

\subsection*{2.2 Likelihood gradient: $M(p)$ and $L(p)$}

For
\[
  M(p) = \frac{\delta(p)^2}{V(p)},
\]
the gradient (using the quotient rule) is
\[
  \nabla M(p)
  =
  \frac{2\delta}{V}\,\nabla\delta
  - \frac{\delta^2}{V^2}\,\nabla V,
\]
so
\[
  \nabla M(p)
  =
  \frac{2\delta}{V}(-N_{\text{ch}}\gamma^\top)
  - \frac{\delta^2}{V^2}N_{\text{ch}}(\sigma^2)^\top
  =
  -N_{\text{ch}}
  \left[
    \frac{2\delta}{V}\,\gamma^\top
    + \frac{\delta^2}{V^2}\,(\sigma^2)^\top
  \right].
\]

For
\[
  L(p) = \log V(p),
\]
\[
  \nabla L(p)
  = \frac{1}{V}\,\nabla V
  = \frac{N_{\text{ch}}}{V}\,(\sigma^2)^\top.
\]

Thus the \emph{exact} gradient of $E_{\text{exact}} = L + M + P$ is
\[
  \begin{aligned}
    \mathbf{g}_{\text{exact}}(p)
    &= \nabla E_{\text{exact}}(p)
     = \nabla P(p) + \nabla M(p) + \nabla L(p)\\[3pt]
    &=
    2N_{\text{ch}}(p-\mu)\Sigma^{-1}
    - N_{\text{ch}}
      \left[
        \frac{2\delta}{V}\,\gamma^\top
        + \frac{\delta^2}{V^2}\,(\sigma^2)^\top
      \right]
    + \frac{N_{\text{ch}}}{V}\,(\sigma^2)^\top.
  \end{aligned}
\]

Grouping $(\sigma^2)^\top$ terms:
\[
  \boxed{
  \mathbf{g}_{\text{exact}}(p)
  =
  2N_{\text{ch}}(p-\mu)\Sigma^{-1}
  -
  N_{\text{ch}}
  \left[
    \frac{2\delta(p)}{V(p)}\,\gamma^\top
    +
    \left(\frac{\delta(p)^2}{V(p)^2} - \frac{1}{V(p)}\right)(\sigma^2)^\top
  \right].
  }
\]

\subsection*{2.3 Likelihood Hessian: $M(p)$ and $L(p)$}

The Hessian of $M(p)$ was derived in detail in the rank-1 setting and
does not change:
\[
  \nabla^2 M(p)
  =
  2N_{\text{ch}}^2
  \left[
    \frac{1}{V}\gamma\gamma^\top
    + \frac{\delta}{V^2}\,
      \bigl(\gamma(\sigma^2)^\top + \sigma^2\gamma^\top\bigr)
    + \frac{\delta^2}{V^3}\,\sigma^2\sigma^2{}^\top
  \right].
\]

For $L(p) = \log V(p)$,
\[
  \nabla^2 L(p)
  = -\frac{1}{V^2}\,\bigl(\nabla V\bigr)^\top\nabla V
  = -\frac{N_{\text{ch}}^2}{V^2}\,\sigma^2\sigma^2{}^\top.
\]

So
\[
  \nabla^2\bigl(M(p)+L(p)\bigr)
  =
  2N_{\text{ch}}^2
  \left[
    \frac{1}{V}\gamma\gamma^\top
    + \frac{\delta}{V^2}\,
      \bigl(\gamma(\sigma^2)^\top + \sigma^2\gamma^\top\bigr)
    + \left(\frac{\delta^2}{V^3} - \frac{1}{2V^2}\right)
      \sigma^2\sigma^2{}^\top
  \right].
\]

\subsection*{2.4 Full exact Hessian and rank-2 structure}

Adding the prior Hessian:
\[
  \nabla^2 P(p) = 2N_{\text{ch}}\Sigma^{-1},
\]
we obtain
\[
  \mathbf{H}_{\text{exact}}(p)
  := \nabla^2 E_{\text{exact}}(p)
  =
  2N_{\text{ch}}\Sigma^{-1}
  +
  2N_{\text{ch}}^2
  \left[
    \frac{1}{V}\gamma\gamma^\top
    + \frac{\delta}{V^2}\,
      \bigl(\gamma(\sigma^2)^\top + \sigma^2\gamma^\top\bigr)
    + \left(\frac{\delta^2}{V^3} - \frac{1}{2V^2}\right)
      \sigma^2\sigma^2{}^\top
  \right].
\]

It is convenient to express this as a rank-2 update of the per-channel
prior precision.  Define
\[
  v(p) = \gamma + \frac{\delta(p)}{V(p)}\sigma^2.
\]

Set
\[
  u_1(p) = v(p),
  \qquad
  u_2(p) = \sigma^2.
\]

We seek a representation
\[
  \mathbf{H}_{\text{exact}}(p)
  = 2N_{\text{ch}}\bigl[\Sigma^{-1} + U(p)C(p)U(p)^\top\bigr],
\]
with
\[
  U(p) = \bigl[u_1(p), u_2(p)\bigr] \in \mathbb{R}^{d\times 2},
\]
and $C(p)$ a diagonal $2\times 2$ matrix.

From the rank-1 analysis without $\log V$ we know that
\[
  \frac{N_{\text{ch}}}{V}v v^\top
  = \frac{N_{\text{ch}}}{V}u_1 u_1^\top
\]
reconstructs the $\gamma\gamma^\top$, $\gamma(\sigma^2)^\top$ and
$\sigma^2\sigma^2{}^\top$ terms in $M$ alone.  The extra contribution
from $L$ is
\[
  -\frac{N_{\text{ch}}^2}{V^2}\sigma^2\sigma^2{}^\top
  = 2N_{\text{ch}}
    \left(-\frac{N_{\text{ch}}}{2V^2}\right) u_2 u_2^\top.
\]

Thus we can take
\[
  C(p) =
  \begin{pmatrix}
    \alpha(p) & 0\\[3pt]
    0         & \beta(p)
  \end{pmatrix},
  \qquad
  \alpha(p) = \frac{N_{\text{ch}}}{V(p)},
  \qquad
  \beta(p)  = -\frac{N_{\text{ch}}}{2V(p)^2},
\]
so that
\[
  U(p)C(p)U(p)^\top
  = \alpha(p)u_1 u_1^\top + \beta(p)u_2 u_2^\top.
\]

We have therefore:
\[
  \boxed{
  \mathbf{H}_{\text{exact}}(p)
  =
  2N_{\text{ch}}
  \bigl[\Sigma^{-1} + U(p)C(p)U(p)^\top\bigr],
  }
\]
a rank-2 update of the (per-channel) prior precision.

\section*{3. Rank-2 Woodbury, Exact Laplace Covariance and Mean}

We now fix an expansion point $p_0$ and compute the Laplace covariance
and mean using the rank-2 structure.  In practice $p_0$ will be taken
as the current Newton iterate or the final MAP $p^\star$.

\subsection*{3.1 Rank-2 Woodbury inversion}

The Woodbury identity for a rank-$k$ update says:
\[
  (A + U C U^\top)^{-1}
  =
  A^{-1}
  - A^{-1}U\bigl(C^{-1} + U^\top A^{-1}U\bigr)^{-1}U^\top A^{-1},
\]
with $A$ invertible, $U$ of shape $d\times k$, and $C$ of shape
$k\times k$.

Formally, on the intrinsic $(d-1)$-dimensional subspace we put
\[
  A = \Sigma^{-1},
  \qquad
  A^{-1} = \Sigma,
  \qquad
  U = U(p_0),
  \qquad
  C = C(p_0).
\]

Then
\[
  \bigl(\Sigma^{-1} + U C U^\top\bigr)^{-1}
  =
  \Sigma
  - \Sigma U K U^\top \Sigma,
\]
where
\[
  K
  =
  \bigl(C^{-1} + U^\top\Sigma U\bigr)^{-1}
  \in \mathbb{R}^{2\times 2}.
\]

At $p_0$ we have
\[
  \mathbf{H}_{\text{exact}}(p_0)
  =
  2N_{\text{ch}}
  \bigl[\Sigma^{-1} + U(p_0)C(p_0)U(p_0)^\top\bigr],
\]
so
\[
  \mathbf{H}_{\text{exact}}(p_0)^{-1}
  =
  \frac{1}{2N_{\text{ch}}}
  \left[
    \Sigma
    - \Sigma U(p_0)K(p_0)U(p_0)^\top \Sigma
  \right].
\]

The exact \emph{actual} Laplace covariance at $p_0$ is
$2\,\mathbf{H}_{\text{exact}}(p_0)^{-1} = \tfrac{1}{N_{\text{ch}}}
\Sigma_{\text{post}}^{\text{exact}}(p_0)$, where the
\emph{scaled} Laplace covariance is
\[
  \boxed{
    \Sigma_{\text{post}}^{\text{exact}}(p_0)
    =
    \Sigma
    - \Sigma\,U(p_0)\,K(p_0)\,U(p_0)^\top\Sigma.
  }
\]
The actual posterior covariance is $\tfrac{1}{N_{\text{ch}}}
\Sigma_{\text{post}}^{\text{exact}}(p_0)$ and shrinks as $1/N_{\text{ch}}$
as expected.  $\Sigma_{\text{post}}^{\text{exact}}$ itself is
$N_{\text{ch}}$-independent in structure.

Explicitly,
\[
  U(p_0) = \bigl[v(p_0),\ \sigma^2\bigr],
  \qquad
  C(p_0) =
  \begin{pmatrix}
    \alpha & 0\\[3pt]
    0      & \beta
  \end{pmatrix},
\]
with
\[
  \alpha = \frac{N_{\text{ch}}}{V(p_0)},
  \qquad
  \beta  = -\frac{N_{\text{ch}}}{2V(p_0)^2},
\]
and
\[
  K(p_0)
  =
  \Bigl(C(p_0)^{-1} + U(p_0)^\top\Sigma U(p_0)\Bigr)^{-1}
  \in \mathbb{R}^{2\times 2}.
\]

The $2\times 2$ matrix
\[
  M = C^{-1} + U^\top\Sigma U
\]
can be written in terms of scalar contractions
\[
  s = v^\top\Sigma v,
  \qquad
  b = v^\top\Sigma\sigma^2,
  \qquad
  c = \sigma^2{}^\top\Sigma\sigma^2,
\]
as
\[
  M =
  \begin{pmatrix}
    \alpha^{-1} + s & b\\[3pt]
    b & \beta^{-1} + c
  \end{pmatrix}
  =
  \begin{pmatrix}
    \dfrac{V(p_0)}{N_{\text{ch}}} + s & b\\[10pt]
    b & -\dfrac{2V(p_0)^2}{N_{\text{ch}}} + c
  \end{pmatrix}.
\]

Then $K(p_0) = M^{-1}$ is obtained with the usual $2\times 2$
inverse, either analytically or numerically.

\paragraph{Remark (singular prior).}
Again, $\Sigma^{-1}$ appears only formally within the Woodbury
identity.  The implementable scaled covariance
$\Sigma_{\text{post}}^{\text{exact}}(p_0)$ contains only $\Sigma$,
$U(p_0)$ and the small $2\times 2$ inverse $K(p_0)$.

\subsection*{3.2 Exact Laplace mean and elimination of $\Sigma^{-1}$}

The Laplace mean at $p_0$ is
\[
  p_{\text{post}}^{\text{exact}}
  =
  p_0 - \mathbf{g}_{\text{exact}}(p_0)\,
         \mathbf{H}_{\text{exact}}(p_0)^{-1}
  =
  p_0 - \frac{1}{2}\mathbf{g}_{\text{exact}}(p_0)\,
         \bigl(\tfrac{1}{N_{\text{ch}}}\Sigma_{\text{post}}^{\text{exact}}(p_0)\bigr),
\]
with $\mathbf{g}_{\text{exact}}$ from Section~2.2.

Introduce
\[
  \Delta p = p_0 - \mu
\]
and write
\[
  \mathbf{g}_{\text{exact}}(p_0)
  =
  2N_{\text{ch}}\Delta p\,\Sigma^{-1}
  - N_{\text{ch}}\,q^\top,
\]
where
\[
  q^\top
  =
  \frac{2\delta}{V}\,\gamma^\top
  +
  \left(\frac{\delta^2}{V^2} - \frac{1}{V}\right)(\sigma^2)^\top,
\]
with $\delta = \delta(p_0)$ and $V = V(p_0)$.

Then
\[
  \begin{aligned}
    p_{\text{post}}^{\text{exact}}
    &=
    p_0
    - \frac{1}{2}\bigl(2N_{\text{ch}}\Delta p\,\Sigma^{-1} - N_{\text{ch}}q^\top\bigr)
      \,\bigl(\tfrac{1}{N_{\text{ch}}}\Sigma_{\text{post}}^{\text{exact}}(p_0)\bigr)\\[3pt]
    &=
    p_0
    - \Delta p\,\Sigma^{-1}\Sigma_{\text{post}}^{\text{exact}}(p_0)
    + \frac{1}{2}\,q^\top
      \Sigma_{\text{post}}^{\text{exact}}(p_0).
  \end{aligned}
\]

Note that the $N_{\text{ch}}$ in the gradient cancels with the
$1/N_{\text{ch}}$ in the actual posterior covariance; this is the
characteristic simplification of the Moffatt scaling.

Using
\[
  \Sigma_{\text{post}}^{\text{exact}}(p_0)
  =
  \Sigma - \Sigma U K U^\top\Sigma,
\]
we obtain
\[
  \begin{aligned}
    \Sigma^{-1}\Sigma_{\text{post}}^{\text{exact}}(p_0)
    &=
    \Sigma^{-1}\bigl(\Sigma - \Sigma U K U^\top\Sigma\bigr)\\[3pt]
    &=
    I - U K U^\top\Sigma,
  \end{aligned}
\]
so
\[
  \Delta p\,\Sigma^{-1}\Sigma_{\text{post}}^{\text{exact}}(p_0)
  =
  \Delta p - \Delta p\,U K U^\top\Sigma.
\]

Therefore
\[
  \begin{aligned}
    p_{\text{post}}^{\text{exact}}
    &=
    p_0
    - \Bigl[\Delta p - \Delta p\,U K U^\top\Sigma\Bigr]
    + \frac{1}{2}q^\top
      \Sigma_{\text{post}}^{\text{exact}}(p_0)\\[3pt]
    &=
    (p_0 - \Delta p)
    + \Delta p\,U K U^\top\Sigma
    + \frac{1}{2}q^\top
      \Sigma_{\text{post}}^{\text{exact}}(p_0).
  \end{aligned}
\]

Since $p_0 - \Delta p = \mu$, we arrive at the exact mean update with
no explicit $\Sigma^{-1}$ and no surviving $N_{\text{ch}}$ factor:
\[
  \boxed{
  p_{\text{post}}^{\text{exact}}
  =
  \mu
  + \Delta p\,U(p_0)K(p_0)U(p_0)^\top\Sigma
  + \frac{1}{2}\,q^\top
    \Sigma_{\text{post}}^{\text{exact}}(p_0),
  }
\]
with
\[
  \Delta p = p_0 - \mu,
  \quad
  q^\top
  =
  \frac{2\delta}{V}\,\gamma^\top
  +
  \left(\frac{\delta^2}{V^2} - \frac{1}{V}\right)(\sigma^2)^\top,
\]
and $\Sigma_{\text{post}}^{\text{exact}}(p_0)$ given above.

The Moffatt scaling makes this update intuitive: the mean depends only
on the per-channel scaled covariance $\Sigma$, and is independent of
$N_{\text{ch}}$ except through the residual $\delta(p_0)$ and the
variance $V(p_0)$, which carry the actual data scale.

\subsection*{3.3 Newton--MAP scheme and MAP--Laplace posterior}

The exact MAP is defined as
\[
  p^\star = \arg\min_{p\in\mathcal{S}} F_{\text{exact}}(p)
  = \arg\min_{p\in\mathcal{S}} E_{\text{exact}}(p),
\]
subject to the simplex constraints.  We locate $p^\star$ via Newton's
method applied to $E_{\text{exact}}$.

Let $p^{(k)}$ be the current iterate.  One Newton step is
\[
  p^{(k+1)} = p^{(k)} - \Delta p^{(k)},
  \qquad
  \mathbf{H}_{\text{exact}}(p^{(k)})\,\Delta p^{(k)\top}
  = \mathbf{g}_{\text{exact}}(p^{(k)})^\top.
\]

Algorithmically:
\begin{enumerate}
  \item Initialise with a feasible point, typically $p^{(0)} = \mu$.
  \item For $k=0,1,2,\dots$ until convergence:
  \begin{enumerate}
    \item Evaluate
      $\delta^{(k)} = \delta(p^{(k)})$,
      $V^{(k)} = V(p^{(k)})$,
      $v^{(k)} = v(p^{(k)})$.
    \item Form
      $\mathbf{g}_{\text{exact}}(p^{(k)})$ and
      $\mathbf{H}_{\text{exact}}(p^{(k)})$ using the formulas above.
    \item Compute the exact scaled Laplace covariance
      $\Sigma_{\text{post}}^{\text{exact}}(p^{(k)})$
      using the rank-2 Woodbury formula.
    \item Update the mean using the exact mean formula with
      $p_0 = p^{(k)}$:
      \[
        p^{(k+1)}
        =
        \mu
        + \Delta p^{(k)} U^{(k)}K^{(k)}U^{(k)\top}\Sigma
        + \frac{1}{2}q^{(k)\top}
          \Sigma_{\text{post}}^{\text{exact}}(p^{(k)}),
      \]
      where $\Delta p^{(k)} = p^{(k)} - \mu$,
      $U^{(k)} = U(p^{(k)})$ and
      $q^{(k)}$ is defined by $\delta^{(k)}$, $V^{(k)}$.
    \item Optionally project $p^{(k+1)}$ back to the simplex
      (enforce $p_i^{(k+1)}\ge 0$ and renormalise).
  \end{enumerate}
\end{enumerate}

After convergence, $p^{(k)}\to p^\star$ and we define the
MAP--Laplace posterior as
\[
  p\mid y
  \approx
  \mathcal{N}\!\bigl(p^\star,\ \tfrac{1}{N_{\text{ch}}}\Sigma_{\text{Lap}}\bigr),
\]
with the \emph{scaled} Laplace covariance
\[
  \Sigma_{\text{Lap}}
  =
  \Sigma_{\text{post}}^{\text{exact}}(p^\star)
  =
  \Sigma - \Sigma\,U(p^\star)\,K(p^\star)\,U(p^\star)^\top\Sigma.
\]

The actual posterior covariance is
$\tfrac{1}{N_{\text{ch}}}\Sigma_{\text{Lap}}$ and inherits the
$1/N_{\text{ch}}$ shrinkage of the Moffatt convention.

\section*{4. Rank-1 Quasi-Laplace Scheme as a Special Case}

We now present the \emph{quasi} (rank-1) scheme obtained by \emph{dropping}
the $\log V(p)$ term from the likelihood.  This corresponds to a
quasi-likelihood / weighted least-squares objective.  It is useful
pedagogically and connects to EKF-like approximations.

\subsection*{4.1 Quasi energy and Hessian}

Define the quasi energy
\[
  E_{\text{quasi}}(p)
  =
  M(p) + P(p)
  =
  \frac{\delta(p)^2}{V(p)}
  + N_{\text{ch}}\,(p-\mu)\,\Sigma^{-1}\,(p-\mu)^\top.
\]

The gradient is
\[
  \mathbf{g}_{\text{quasi}}(p)
  =
  2N_{\text{ch}}(p-\mu)\Sigma^{-1}
  -
  N_{\text{ch}}
  \left[
    \frac{2\delta(p)}{V(p)}\,\gamma^\top
    + \frac{\delta(p)^2}{V(p)^2}\,(\sigma^2)^\top
  \right].
\]

The Hessian is
\[
  \mathbf{H}_{\text{quasi}}(p)
  =
  2N_{\text{ch}}\Sigma^{-1}
  +
  2N_{\text{ch}}^2
  \left[
    \frac{1}{V}\gamma\gamma^\top
    + \frac{\delta}{V^2}\,
      \bigl(\gamma(\sigma^2)^\top + \sigma^2\gamma^\top\bigr)
    + \frac{\delta^2}{V^3}\,\sigma^2\sigma^2{}^\top
  \right].
\]

Introduce
\[
  v(p) = \gamma + \frac{\delta(p)}{V(p)}\sigma^2,
  \qquad
  u(p) = \sqrt{\frac{N_{\text{ch}}}{V(p)}}\,v(p).
\]

Then
\[
  \mathbf{H}_{\text{quasi}}(p)
  =
  2N_{\text{ch}}
  \left[
    \Sigma^{-1}
    + u(p)u(p)^\top
  \right].
\]

This is a rank-1 update of the per-channel prior precision.

\subsection*{4.2 Rank-1 scaled covariance and mean}

Applying Sherman--Morrison with $A=\Sigma^{-1}$ and
$u = u(p_0)$ at an expansion point $p_0$:
\[
  \bigl(\Sigma^{-1} + u u^\top\bigr)^{-1}
  =
  \Sigma
  - \frac{\Sigma u u^\top\Sigma}{1 + u^\top\Sigma u}.
\]

Therefore
\[
  \mathbf{H}_{\text{quasi}}(p_0)^{-1}
  =
  \frac{1}{2N_{\text{ch}}}
  \left[
    \Sigma
    - \frac{\Sigma u u^\top\Sigma}{1 + u^\top\Sigma u}
  \right].
\]

The actual quasi Laplace covariance at $p_0$ is
$2\,\mathbf{H}_{\text{quasi}}(p_0)^{-1} = \tfrac{1}{N_{\text{ch}}}
\Sigma_{\text{post}}^{\text{quasi}}(p_0)$, and the
\emph{scaled} quasi Laplace covariance is
\[
  \Sigma_{\text{post}}^{\text{quasi}}(p_0)
  =
  \Sigma
  - \frac{\Sigma u u^\top\Sigma}{1 + u^\top\Sigma u}.
\]

Substituting $u = \sqrt{N_{\text{ch}}/V}\,v$ and defining
$s = v^\top\Sigma v$, we obtain
\[
  \boxed{
    \Sigma_{\text{post}}^{\text{quasi}}(p_0)
    =
    \Sigma
    -
    \frac{N_{\text{ch}}}{V(p_0) + N_{\text{ch}} s(p_0)}\,
    \Sigma v(p_0) v(p_0)^\top\Sigma.
  }
\]

This matches Eq.~(15) of the gemmini-Moffatt note exactly.

The quasi mean at $p_0$ is
\[
  p_{\text{post}}^{\text{quasi}}
  =
  p_0 - \frac{1}{2}\mathbf{g}_{\text{quasi}}(p_0)\,
         \bigl(\tfrac{1}{N_{\text{ch}}}\Sigma_{\text{post}}^{\text{quasi}}(p_0)\bigr).
\]

Writing $\Delta p = p_0 - \mu$ and using the identity
\[
  \Sigma^{-1}\Sigma_{\text{post}}^{\text{quasi}}
  =
  I
  - \frac{N_{\text{ch}}}{V + N_{\text{ch}} s}\,v v^\top\Sigma,
\]
we obtain, after the same algebra as in Section~3.2 (with the
$N_{\text{ch}}$ cancelling against $1/N_{\text{ch}}$):
\[
  \boxed{
  \begin{aligned}
    p_{\text{post}}^{\text{quasi}}
    &=
    \mu
    + \frac{N_{\text{ch}}}{V(p_0) + N_{\text{ch}} s(p_0)}\,
      (p_0 - \mu)\,v(p_0)v(p_0)^\top\Sigma\\[4pt]
    &\quad
    + \frac{\delta(p_0)}{V(p_0)}\,\gamma^\top
      \Sigma_{\text{post}}^{\text{quasi}}(p_0)
    + \frac{\delta(p_0)^2}{2V(p_0)^2}\,(\sigma^2)^\top
      \Sigma_{\text{post}}^{\text{quasi}}(p_0).
  \end{aligned}
  }
\]

Again, no explicit $\Sigma^{-1}$ survives and no $N_{\text{ch}}$ factor
multiplies the scaled posterior covariance.

\subsection*{4.3 Special case $p_0 = \mu$ and connection to EKF}

If we choose $p_0 = \mu$, then $\Delta p = 0$ and the first line of
the quasi mean vanishes.  In that case:
\[
  \mathbf{g}_{\text{quasi}}(\mu)
  =
  -N_{\text{ch}}
  \left[
    \frac{2\delta}{V}\,\gamma^\top
    + \frac{\delta^2}{V^2}\,(\sigma^2)^\top
  \right]
  =
  -\frac{N_{\text{ch}}\delta}{V}\bigl(\gamma^\top + v^\top\bigr),
\]
and
\[
  p_{\text{post}}^{\text{quasi}}
  =
  \mu + \frac{\delta}{2V}\,(\gamma^\top + v^\top)
        \Sigma_{\text{post}}^{\text{quasi}},
\]
where $v = \gamma + \frac{\delta}{V}\sigma^2$ and
$\Sigma_{\text{post}}^{\text{quasi}}$ is evaluated at $p_0 = \mu$.

This is the form quoted in the deepseek-Moffatt note:
\[
  r_{\text{post}}
  =
  \mu
  + \frac{\delta}{2V}\,\Sigma_{\text{post}}\,(\gamma + v).
\]
It can be interpreted as a single quasi-Laplace (or EKF-like) update
from the prior mean, using only the mean-matching part of the
likelihood and ignoring the $\log V$ term.

\subsection*{4.4 Recovery of Moffatt (2007): $\sigma^2 = 0$}

When state-dependent intrinsic noise vanishes, $\sigma^2 = 0$,
the tilted vector simplifies:
\[
  v = \gamma,
  \qquad
  V = \epsilon^2,
\]
and the quasi-Laplace formulas reduce to
\[
  \Sigma_{\text{post}}^{\text{quasi}}(p_0)
  =
  \Sigma
  - \frac{N_{\text{ch}}}{\epsilon^2 + N_{\text{ch}}\,\gamma^\top\Sigma\gamma}\,
    \Sigma\gamma\gamma^\top\Sigma,
\]
\[
  p_{\text{post}}^{\text{quasi}}\big|_{p_0=\mu}
  =
  \mu
  + \frac{\delta}{\epsilon^2}\,\gamma^\top\Sigma_{\text{post}}^{\text{quasi}}.
\]
These are exactly the single-step measurement-update equations of
Moffatt (2007) (Eqs.~49--50 in the gemmini-Moffatt note's numbering).
The Moffatt scaling makes this recovery transparent: $\Sigma$,
$\Sigma_{\text{post}}$ and the mean update all reduce term-by-term to
the 2007 expressions, with the $1/N_{\text{ch}}$ shrinkage living
exclusively in the actual covariance of the Gaussian.

\section*{5. Summary and Interpretation}

We now summarise the two schemes under Moffatt scaling.

\subsection*{5.1 Rank-2 (exact) scheme}

\begin{itemize}
  \item Likelihood:
    \[
      p(y\mid p) =
      \mathcal{N}\bigl(N_{\text{ch}}(p\cdot\gamma), V(p)\bigr),
      \qquad
      V(p) = \epsilon^2 + N_{\text{ch}}(p\cdot\sigma^2).
    \]
  \item Prior (Moffatt-scaled):
    \[
      p \sim \mathcal{N}\!\bigl(\mu,\ \tfrac{1}{N_{\text{ch}}}\Sigma\bigr),
      \qquad
      \Lambda_p = N_{\text{ch}}\Sigma^{-1}.
    \]
  \item Negative log-posterior (up to const):
    \[
      F_{\text{exact}}(p)
      =
      \frac{1}{2}\log V(p)
      + \frac{1}{2}\frac{\delta(p)^2}{V(p)}
      + \frac{N_{\text{ch}}}{2}(p-\mu)\,\Sigma^{-1}\,(p-\mu)^\top.
    \]
  \item Energy:
    \[
      E_{\text{exact}}(p) = 2F_{\text{exact}}(p).
    \]
  \item Hessian:
    \[
      \mathbf{H}_{\text{exact}}(p)
      =
      2N_{\text{ch}}
      \bigl[\Sigma^{-1} + U(p)C(p)U(p)^\top\bigr],
    \]
    with $U(p) = [v(p),\sigma^2]$ and diagonal
    $C(p) = \mathrm{diag}\bigl(N_{\text{ch}}/V,\ -N_{\text{ch}}/(2V^2)\bigr)$.
  \item Scaled covariance:
    \[
      \Sigma_{\text{post}}^{\text{exact}}(p_0)
      =
      \Sigma
      - \Sigma\,U(p_0)\,K(p_0)\,U(p_0)^\top\Sigma,
    \]
    with $K(p_0) = \bigl(C(p_0)^{-1} + U(p_0)^\top\Sigma U(p_0)\bigr)^{-1}$.
    The actual posterior covariance is
    $\tfrac{1}{N_{\text{ch}}}\Sigma_{\text{post}}^{\text{exact}}(p_0)$.
  \item Mean:
    \[
      p_{\text{post}}^{\text{exact}}
      =
      \mu
      + \Delta p\,U(p_0)K(p_0)U(p_0)^\top\Sigma
      + \frac{1}{2}q^\top
        \Sigma_{\text{post}}^{\text{exact}}(p_0).
    \]
  \item Newton iteration uses $E_{\text{exact}}$, converges to the
        true MAP $p^\star$, and the scaled Laplace covariance at
        $p^\star$ is
        $\Sigma_{\text{Lap}} = \Sigma_{\text{post}}^{\text{exact}}(p^\star)$.
\end{itemize}

This scheme fully incorporates the heteroscedastic variance $V(p)$ via
both the $\delta^2/V$ term and the $\log V$ term of the Gaussian
likelihood.  Even when $\delta=0$, the $\log V$ term can shift the MAP
towards states with smaller predicted variance, in competition with
the prior.

\subsection*{5.2 Rank-1 (quasi) scheme}

\begin{itemize}
  \item Energy:
    \[
      E_{\text{quasi}}(p)
      =
      \frac{\delta(p)^2}{V(p)}
      + N_{\text{ch}}\,(p-\mu)\,\Sigma^{-1}\,(p-\mu)^\top.
    \]
  \item Hessian:
    \[
      \mathbf{H}_{\text{quasi}}(p)
      =
      2N_{\text{ch}}
      \left[
        \Sigma^{-1}
        + \frac{N_{\text{ch}}}{V(p)}\,v(p)v(p)^\top
      \right],
    \]
    a rank-1 update.
  \item Scaled covariance:
    \[
      \Sigma_{\text{post}}^{\text{quasi}}(p_0)
      =
      \Sigma
      -
      \frac{N_{\text{ch}}}{V(p_0) + N_{\text{ch}} s(p_0)}\,
      \Sigma v(p_0)v(p_0)^\top\Sigma.
    \]
  \item Mean: the quasi mean $p_{\text{post}}^{\text{quasi}}$ is given
    by the rank-1 formula above, and in the special case $p_0=\mu$
    reduces to a single-step update
    \[
      p_{\text{post}}^{\text{quasi}}
      =
      \mu
      + \frac{\delta}{2V}\,(\gamma^\top + v^\top)
        \Sigma_{\text{post}}^{\text{quasi}}.
    \]
\end{itemize}

This scheme ignores the $\log V(p)$ contribution to the likelihood and
treats $V(p)$ only as a weight in the squared residual.  It is
analogous to using a weighted least-squares / quasi-likelihood
objective in GLMs, and connects to EKF practice where a
state-dependent variance is plugged into an effective $R$ but not
differentiated.  When $\sigma^2 = 0$ it reduces to Moffatt (2007).

\subsection*{5.3 Practical recommendation}

In MacroTaylor we take the rank-2 scheme as the \emph{main method}:
\begin{itemize}
  \item it is the correct Gaussian MAP for the interval model;
  \item the rank-2 Woodbury cost is negligible relative to the cost of
        computing the MacroIR quantities;
  \item the heteroscedastic variance $V(p)$ is physically meaningful
        and carries information about $p$.
\end{itemize}

The rank-1 scheme is retained as a \emph{pedagogical device} and as a
limiting case:
\begin{itemize}
  \item it illustrates how a simpler quasi-likelihood/EKF update
        arises if one drops $\log V(p)$;
  \item it provides a useful check in regimes where $V(p)$ changes
        slowly with $p$ or $N_{\text{ch}}$ is large;
  \item with $\sigma^2 = 0$ it recovers Moffatt (2007) directly.
\end{itemize}

\subsection*{5.4 What the Moffatt scaling buys}

Comparing this note with \texttt{MacroTaylor\_2\_chatgpt\_reseach.tex}
(which used the unscaled convention $p \sim \mathcal{N}(\mu,\Sigma_p)$
with $\Lambda_p = N_{\text{ch}}\Sigma_p^{-1}$):

\begin{itemize}
  \item The scaled prior covariance $\Sigma$ here is
        $N_{\text{ch}}$-independent; in the unscaled version,
        $\Sigma_p$ played the same role but was conceptually the full
        population covariance.  Numerically they coincide.
  \item The scaled posterior covariance $\Sigma_{\text{post}}$ now has
        \emph{no} explicit $1/N_{\text{ch}}$ factor in front; the
        $1/N_{\text{ch}}$ shrinkage lives in the Gaussian
        $\mathcal{N}\bigl(p^\star, \tfrac{1}{N_{\text{ch}}}\Sigma_{\text{post}}\bigr)$
        rather than in $\Sigma_{\text{post}}$ itself.
  \item The mean update has no $N_{\text{ch}}/2$ factor on the
        $q^\top\Sigma_{\text{post}}$ term: the $N_{\text{ch}}$ in the
        gradient cancels with the $1/N_{\text{ch}}$ from the actual
        posterior covariance.  Compare
        \[
          p_{\text{post}}^{\text{exact}}
          =
          \mu + \Delta p\,UKU^\top\Sigma + \tfrac{1}{2}\,q^\top\Sigma_{\text{post}}^{\text{exact}}
        \]
        (Moffatt scaling, this note) with
        \[
          p_{\text{post}}^{\text{exact}}
          =
          \mu + \Delta p\,UKU^\top\Sigma_p + \tfrac{N_{\text{ch}}}{2}\,q^\top\Sigma_{\text{post}}^{\text{exact}}
        \]
        (unscaled, original \texttt{MacroTaylor\_2}); the two are
        algebraically equivalent under
        $\Sigma \leftrightarrow \Sigma_p$ and
        $\Sigma_{\text{post}} \leftrightarrow N_{\text{ch}}\Sigma_{\text{post}}^{\text{tex}}$.
  \item The limit $\sigma^2 \to 0$ reduces directly to the single-step
        update of Moffatt (2007), without any rescaling of either
        $\Sigma$ or $\Sigma_{\text{post}}$.
\end{itemize}

The Moffatt scaling is therefore the natural choice when (i) comparing
to the 2007 single-channel formulas, (ii) reasoning about the
$N_{\text{ch}}\to\infty$ thermodynamic limit (in which only
$\tfrac{1}{N_{\text{ch}}}\Sigma_{\text{post}}$ shrinks), and
(iii) keeping all matrices in the algebra $N_{\text{ch}}$-independent
in structure.

\end{document}
